\chapter{Provas sobre transacções}
\label{chapter:tx-knowledge-proofs}
% Original draft by Sarang Noether, Ph.D.; implementation review for this
% edition against Monero revision 3646f648db57f60cca86430e25a635d19fa9b92a.

A lista de blocos não publica os endereços dos remetentes e destinatários, e
os montantes RingCT estão cifrados. Uma carteira pode, porém, divulgar provas
limitadas sobre uma transacção ou sobre um conjunto de saídas. Este capítulo
define as afirmações que essas provas sustentam, os dados que revelam e as
condições que não demonstram.

Salvo indicação em contrário, as observações sobre a implementação referem-se
à revisão do código Monero datada de 14 de Agosto de 2026
\cite{monero-source-v16}. Estas provas são artefactos da carteira. Não fazem
parte de uma transacção, não são verificadas pelo consenso e não alteram o
estado da lista de blocos.

\section{Interface e alcance}
\label{sec:proofs-monero-proofs}

A carteira consultada gera quatro formatos:

\begin{center}
\small
\begin{tabular}{|>{\raggedright\arraybackslash}p{3.25cm}|
  >{\raggedright\arraybackslash}p{3.4cm}|
  >{\raggedright\arraybackslash}p{7.85cm}|}
\hline
\textbf{Formato} & \textbf{Segredo usado} & \textbf{Afirmação verificada} \\
\hline
\texttt{OutProofV2} & chave privada de transacção & O endereço indicado recebeu
um montante determinado na transacção. \\
\hline
\texttt{InProofV2} & chave privada de vista & O endereço indicado, controlado
para efeitos de vista pelo provador, recebeu um montante determinado. \\
\hline
\texttt{SpendProofV1} & chaves privadas das saídas gastas & O provador conhece
a chave de gasto correspondente a uma posição desconhecida de cada anel de
entrada. \\
\hline
\texttt{ReserveProofV2} & chaves privadas de vista e de gasto & O provador
controla as saídas enumeradas e liga a prova a um endereço principal; o
verificador calcula o total e a parte já gasta. \\
\hline
\end{tabular}
\end{center}

Os verificadores ainda aceitam \texttt{InProofV1}, \texttt{OutProofV1} e
\texttt{ReserveProofV1} para compatibilidade com artefactos antigos. A carteira
gera as versões V2. A prova de gasto continua a usar
\texttt{SpendProofV1}; não existe no código consultado uma
\texttt{SpendProofV2}.

Na tabela, ``recebeu'' não significa ``permanece não gasta''. Uma prova de
entrada ou de saída confirma a recepção e o montante, mas não determina se a
saída foi gasta depois, não reconstrói a actividade de saída da carteira e
não prova o saldo actual.

Na carteira de linha de comandos, as operações agrupam-se assim:
\begin{itemize}
 \item \texttt{get\_tx\_proof} e \texttt{check\_tx\_proof};
 \item \texttt{get\_spend\_proof} e \texttt{check\_spend\_proof};
 \item \texttt{get\_reserve\_proof} e \texttt{check\_reserve\_proof}.
\end{itemize}
A interface RPC da carteira expõe métodos com os mesmos nomes. A geração e a
verificação recebem uma mensagem opcional. As duas partes têm de usar
exactamente os mesmos bytes.

Para as provas relativas a uma só transacção, a mensagem assinada é
\begin{align*}
 \mathfrak m=\mathcal H(\texttt{txid}\mathbin\|\texttt{message}).
\end{align*}
O campo \texttt{message} deve conter o contexto da afirmação, por exemplo um
identificador de factura e um desafio aleatório escolhido pelo verificador.
Uma mensagem vazia produz uma prova transferível: qualquer pessoa que a
receba pode apresentá-la de novo.

Os prefixos de formato ficam em texto simples. O corpo é serializado e
codificado na variante Base58 de Monero \cite{base-58-encoding}. Para verificar
uma prova são ainda necessários o ID da transacção, o endereço e a mensagem
correctos, bem como acesso a um nó que forneça a transacção e as saídas
referidas.

\subsection{A prova V2 com duas bases}
\label{subsec:proofs-multi-base-monero}

\texttt{OutProofV2}, \texttt{InProofV2} e uma parte de
\texttt{ReserveProofV2} usam a mesma prova de igualdade de logaritmos
discretos. Sejam
\begin{align*}
 P_1&=xJ_1, & P_2&=xJ_2,
\end{align*}
onde o provador conhece \(x\). Definimos \(J_1^*=0^{32}\) quando
\(J_1=G\), e \(J_1^*=J_1\) quando \(J_1\neq G\). Esta convenção reproduz o
último campo do traslado implementado. A separação de domínio é
\begin{align*}
 T_{\rm tx}=\mathcal H(\texttt{TXPROOF\_V2}).
\end{align*}

O provador escolhe \(\alpha\in_R\mathbb Z_l\), calcula
\begin{align*}
 L_1&=\alpha J_1, & L_2&=\alpha J_2,\\
 c&=\mathcal H_n(\mathfrak m,P_2,L_1,L_2,T_{\rm tx},P_1,J_2,J_1^*),\\
 z&=\alpha-cx\pmod l,
\end{align*}
e publica \((c,z)\). O verificador recompõe
\begin{align*}
 L'_1&=zJ_1+cP_1, & L'_2&=zJ_2+cP_2
\end{align*}
e aceita quando o desafio recalculado com \(L'_1,L'_2\) é igual a \(c\).
Assim, a prova liga o mesmo escalar às duas relações sem o revelar.

\section{Revelar a chave privada de transacção}
\label{sec:proofs-transaction-key}

Antes das provas codificadas, há uma opção mais simples. A carteira remetente
pode obter a chave privada de transacção com \texttt{get\_tx\_key}; o
verificador usa \texttt{check\_tx\_key} com o ID, a chave e o endereço. Para um
endereço principal, a relação pública é
\begin{align*}
 R=rG,
\end{align*}
e o verificador calcula a derivação a partir de \(rK^v\). Para uma transacção
com chaves adicionais, o valor devolvido concatena a chave principal e todas
as chaves adicionais. O verificador testa as derivações contra as saídas,
descodifica os montantes que pertencem ao endereço e devolve o total recebido,
o estado na memória temporária e o número de confirmações.

A chave de transacção não permite gastar as saídas nem obter as chaves privadas
do destinatário. Contudo, permite testar endereços candidatos e identificar
os que receberam fundos nessa transacção. Também permite produzir uma
\texttt{OutProofV2}. Deve, por isso, ser entregue apenas quando essa divulgação
for aceitável. Se a carteira remetente não conservou as chaves, não pode
recuperá-las da lista de blocos.

\section{Prova de saída: \texttt{OutProofV2}}
\label{subsec:proofs-output-creator-outproof}

Uma prova de saída demonstra a relação entre a chave pública de transacção, o
endereço indicado e o segredo partilhado, sem revelar a chave privada de
transacção. Para um pagamento a um endereço principal \((K^s,K^v)\), a prova
de duas bases usa
\begin{align*}
 x&=r, & J_1&=G, & J_2&=K^v,\\
 P_1&=R=rG, & P_2&=D=rK^v.
\end{align*}
Para um sub-endereço \((K^{s,i},K^{v,i})\), usa
\begin{align*}
 x&=r, & J_1&=K^{s,i}, & J_2&=K^{v,i},\\
 P_1&=R=rK^{s,i}, & P_2&=D=rK^{v,i}.
\end{align*}

O corpo de \texttt{OutProofV2} contém, para a chave pública principal e para
cada chave pública adicional da transacção, o ponto partilhado \(D\) e a
assinatura \((c,z)\). O verificador confirma as assinaturas, converte cada
\(D\) na derivação \(8D\), testa as chaves de saída contra o endereço e
descodifica os respectivos montantes RingCT. A prova só é gerada se a pesquisa
encontrar um montante não nulo para o endereço apresentado.

O resultado prova que quem construiu o artefacto conhecia a chave privada de
transacção correspondente e que o endereço recebeu o montante calculado. Não
prova a identidade civil do remetente, não identifica as entradas usadas e
não prova que o provador foi o único participante na construção da
transacção. Quem obtiver a chave privada de transacção também consegue gerar a
mesma classe de prova.

\section{Prova de entrada: \texttt{InProofV2}}
\label{subsec:proofs-output-ownership-inproof}

Quando o endereço apresentado pertence à carteira, \texttt{get\_tx\_proof}
gera uma prova de entrada. Para um endereço principal, sejam \(R=rG\) a chave
pública de transacção e \(k^v\) a chave privada de vista. A prova usa
\begin{align*}
 x&=k^v, & J_1&=G, & J_2&=R,\\
 P_1&=K^v=k^vG, & P_2&=D=k^vR.
\end{align*}
Para um sub-endereço, substitui-se a primeira base por \(K^{s,i}\) e a
primeira chave pública por
\(K^{v,i}=k^vK^{s,i}\); a segunda relação continua a ser
\(D=k^vR\).

O formato serializado é paralelo ao de \texttt{OutProofV2}: o prefixo
\texttt{InProofV2} é seguido por um segredo partilhado e uma assinatura para
cada chave pública de transacção. O verificador confirma a igualdade das duas
relações, deriva \(8D\), encontra as saídas dirigidas ao endereço e soma os
montantes.

Esta prova demonstra conhecimento da chave privada de vista associada ao
endereço e a recepção do montante indicado. Não demonstra conhecimento da
chave privada de gasto e, portanto, não prova por si só que o provador possa
gastar a saída. Uma carteira só de vista pode gerar esta classe de prova. A
divulgação de provas de entrada escolhidas não entrega \(k^v\); a divulgação
da própria chave de vista permitiria ao verificador detectar também as
recepções futuras da conta.

\subsection{O estado gasto e a actividade de saída}
\label{subsec:proofs-view-key-limit}

Na hierarquia actualmente usada, a chave privada de vista detecta uma saída
recebida e descodifica o montante. Para decidir se essa saída foi gasta, a
carteira tem de procurar na lista de blocos a sua imagem de chave. No caso
principal,
\begin{align*}
 k^o&=\mathcal H_n\bigl(8k^vR,\operatorname{varint}(t)\bigr)+k^s\pmod l,\\
 I&=k^o\mathcal H_p(K^o).
\end{align*}
A chave de vista \(k^v\) não fornece a chave de gasto \(k^s\); por isso, não
fornece \(k^o\) nem a imagem \(I\). A lista de blocos publica \(I\) quando a
saída é gasta, mas não publica uma ligação que a chave de vista consiga
reconstruir sozinha.

Consequentemente, uma carteira restaurada apenas com o endereço e \(k^v\)
consegue somar recepções, mas não consegue, a partir desses dados e da lista
de blocos, subtrair com segurança as saídas gastas, reconstruir as transacções
de saída ou calcular o saldo actual. Pode detectar uma saída de troco como
uma nova recepção, sem com isso identificar o gasto que a produziu nem o seu
destinatário. Se uma carteira só de vista apresentar o estado gasto ou o
histórico de saída, essa informação provém de imagens de chave ou de metadados
fornecidos por uma carteira completa, não da chave de vista isolada.

\texttt{InProofV2} não elimina esta limitação: prova a relação de vista e o
montante recebido, mas não fornece a imagem de chave. Em contraste,
\texttt{ReserveProofV2} é gerada por uma carteira completa, inclui as imagens
de chave das saídas declaradas e permite ao verificador consultar o estado de
gasto dessas imagens.

\subsection{Âmbitos de vista previstos no Carrot}
\label{subsec:proofs-carrot-view-tiers}

O protocolo de endereçamento Carrot, previsto para a actualização FCMP++,
separa capacidades que a hierarquia anterior concentra ou não oferece
\cite{carrot-specification}. Na hierarquia recomendada:
\begin{itemize}
 \item a chave de vista de entrada \(k_v\), correspondente ao nível
 \emph{View-Received}, detecta pagamentos externos e troco externo, mas não
 mostra onde as saídas foram gastas nem as saídas internas;

 \item o segredo de vista de saldo \(s_{vb}\), correspondente ao nível
 \emph{View-All}, permite derivar a informação necessária para reconhecer
 recepções externas e internas, gastos e o saldo, sem autorizar esses gastos;

 \item o segredo \(s_{ga}\) permite gerar endereços sem conceder acesso à
 lista de blocos.
\end{itemize}
Assim, a capacidade de observar envios é fornecida pelo nível associado a
\(s_{vb}\), não pela chave privada de vista legada. Os formatos descritos
neste capítulo são os formatos
da carteira actual. A especificação Carrot ainda não define versões desses
formatos para saídas Carrot nem afirma que possam ser reutilizados sem
adaptação.

\section{Prova de gasto: \texttt{SpendProofV1}}
\label{subsec:proofs-input-creation-spendproof}

Uma prova de gasto contém uma nova assinatura em anel para cada entrada de
gasto; no código, estas entradas têm o tipo \texttt{txin\_to\_key}. A carteira
reutiliza exactamente o anel e a imagem de chave da entrada original.
Reutilizar os mesmos membros evita uma intersecção entre dois anéis diferentes
que pudesse reduzir o anonimato.

Embora as transacções activas usem CLSAG, esta prova auxiliar usa o esquema de
assinatura em anel original de CryptoNote \cite{cryptoNoteWhitePaper}. Para um
anel \((K_1^o,\ldots,K_n^o)\), a posição real \(\pi\), a chave privada
\(k_\pi^o\) e a imagem de chave
\begin{align*}
 \widetilde K=k_\pi^o\mathcal H_p(K_\pi^o),
\end{align*}
o provador escolhe \(\alpha\in_R\mathbb Z_l\) e pares aleatórios
\((c_i,z_i)\) para \(i\neq\pi\). Define
\begin{align*}
 L_i&=z_iG+c_iK_i^o,&
 R_i&=z_i\mathcal H_p(K_i^o)+c_i\widetilde K
 \qquad(i\neq\pi),\\
 L_\pi&=\alpha G,&
 R_\pi&=\alpha\mathcal H_p(K_\pi^o),\\
 c_{\rm tot}&=\mathcal H_n(\mathfrak m,L_1,R_1,\ldots,L_n,R_n),\\
 c_\pi&=c_{\rm tot}-\sum_{i\neq\pi}c_i\pmod l,\\
 z_\pi&=\alpha-c_\pi k_\pi^o\pmod l.
\end{align*}
Publica os pares \((c_1,z_1),\ldots,(c_n,z_n)\). O verificador recompõe todos
os \(L_i,R_i\) e aceita se
\begin{align*}
 \mathcal H_n(\mathfrak m,L_1,R_1,\ldots,L_n,R_n)
 \stackrel{?}{=}\sum_{i=1}^n c_i\pmod l.
\end{align*}

\texttt{SpendProofV1} concatena essas assinaturas codificadas em Base58. A
carteira geradora tem de possuir as chaves correspondentes a todas as imagens
de chave da transacção. Uma carteira só de vista e a interface de uma carteira
de hardware não geram esta prova.

Uma prova válida demonstra o controlo de uma saída real em cada anel de
entrada, sem revelar qual membro era o real. Não identifica o destinatário,
não revela os montantes enviados e não demonstra, isoladamente, a identidade
do autor. Numa transacção construída em colaboração, a ferramenta só funciona
para uma carteira que disponha das chaves de todas as entradas.

\section{Prova de reserva: \texttt{ReserveProofV2}}
\label{subsec:proofs-minimum-balance-reserveproof}

Uma prova de reserva associa um endereço principal a saídas controladas pela
carteira. Pode abranger todas as saídas elegíveis ou apenas as saídas
necessárias para ultrapassar um montante mínimo. No segundo caso, a carteira
selecciona primeiro as saídas de maior valor e pode provar um total superior
ao mínimo pedido. Saídas gastas ou congeladas não entram na selecção.

A geração exige uma carteira completa, não multi-assinatura, com as chaves
privadas de vista e de gasto. A carteira de linha de comandos também recusa a
operação através de uma carteira de hardware. O verificador fornece o endereço
principal; um sub-endereço não é aceite nesse campo.

Sejam \(I_1,\ldots,I_q\) as imagens de chave das saídas seleccionadas. Todas
as componentes assinam
\begin{align*}
 \mathfrak m_{\rm res}=\mathcal H(
 \texttt{message}\mathbin\|\texttt{address}\mathbin\|I_1\mathbin\|\cdots
 \mathbin\|I_q).
\end{align*}
Aqui \texttt{address} é a representação binária do par de chaves públicas,
não o texto Base58. O corpo de \texttt{ReserveProofV2} contém duas colecções:

\begin{enumerate}
 \item Para cada saída: o ID da transacção, o índice da saída, o segredo
 partilhado \(D=k^vR\), a imagem de chave, uma prova V2 do segredo partilhado
 e uma assinatura que liga a imagem de chave à chave de saída.

 \item Para cada chave pública de gasto principal ou de sub-endereço que
 controla alguma das saídas: uma assinatura de Schnorr com a chave privada de
 gasto correspondente.
\end{enumerate}

Na verificação, a carteira rejeita entradas ou imagens de chave duplicadas,
obtém as transacções confirmadas, verifica as provas do segredo partilhado e
das imagens de chave, deriva o sub-endereço pertinente, descodifica cada
montante e verifica as assinaturas das chaves de gasto. Consulta ainda o nó
para classificar cada imagem de chave como gasta ou não gasta. O resultado é
o par de montantes \texttt{total} e \texttt{spent}; a reserva não gasta no
momento da consulta é
\begin{align*}
 \texttt{total}-\texttt{spent}.
\end{align*}

A prova revela as saídas seleccionadas, os seus montantes, as imagens de
chave e as chaves públicas de gasto dos sub-endereços envolvidos. Uma imagem
de chave ainda não vista permite reconhecer publicamente o gasto futuro da
saída. As assinaturas demonstram o controlo separado dessas chaves de gasto,
mas não provam que cada sub-endereço foi derivado canonicamente da chave de
gasto principal.

A prova de um mínimo estabelece apenas um limite inferior; o provador pode
omitir outras contas, outras carteiras e fundos não seleccionados. Mesmo a
opção \texttt{all} cobre apenas as saídas elegíveis conhecidas pela carteira
que produziu o artefacto. O estado gasto ou não gasto deve ser consultado de
novo, porque pode mudar depois da geração.

\section{Afirmações de pagamento e identidade}

As provas anteriores tratam relações entre chaves. Não autenticam nomes,
empresas, facturas ou contractos. Para ligar uma prova a uma identidade, o
verificador deve fornecer uma mensagem específica e exigir, separadamente,
uma assinatura autenticada dessa mensagem. A mensagem deve incluir pelo menos
o ID da transacção, o endereço, a finalidade da prova e um desafio fresco.

As afirmações disponíveis podem resumir-se assim:

\begin{itemize}
 \item uma chave de transacção ou \texttt{OutProofV2} demonstra que um endereço
 recebeu certo montante, mas não identifica quem controlava as entradas;

 \item \texttt{InProofV2} demonstra recepção e controlo da chave de vista, mas
 não autorização de gasto;

 \item \texttt{SpendProofV1} demonstra controlo das entradas, mas não o destino
 nem o montante líquido enviado a uma pessoa;

 \item \texttt{ReserveProofV2} demonstra controlo das saídas enumeradas e o
 seu estado observado, mas não a totalidade do património do provador.
\end{itemize}

Uma auditoria que precise de conhecer a origem declarada de um pagamento deve
combinar estas provas com documentos externos e autenticação das partes. A
lista de blocos e as provas da carteira, por si sós, não fornecem essa
identidade.

\section{Construções não implementadas}

As duas construções seguintes descrevem relações úteis para auditoria, mas
não existem nas interfaces da carteira consultada. Não são formatos de prova
aceites pelas ferramentas actuais.

\subsection{Provar que uma saída não corresponde a uma imagem de chave}
\label{subsec:proofs-owned-output-spent-unspentproof}

Considere-se uma saída própria \(K^o\), com
\begin{align*}
 k^o=h+k^s\pmod l,
 \qquad h=\mathcal H_n\bigl(8rK^v,\operatorname{varint}(t)\bigr),
\end{align*}
e uma imagem de chave \(I_?\) publicada por uma entrada cujo anel contém
\(K^o\). O verificador calcula
\begin{align*}
 Z_?=I_?-h\mathcal H_p(K^o).
\end{align*}
Se essa entrada gastou \(K^o\), então
\begin{align*}
 Z_?=k^s\mathcal H_p(K^o).
\end{align*}

A proposta de prova de não-gasto \cite{unspent-proof-issue-68} evita publicar
directamente esse último ponto. O provador demonstra duas relações:
\begin{enumerate}
 \item com testemunha \(k^s\),
 \begin{align*}
  \mathcal J_3&=\{G,K^s,Z_?\},\\
  \mathcal K_3&=\{K^s,k^sK^s,k^sZ_?\};
 \end{align*}

 \item com testemunha \((k^s)^2\),
 \begin{align*}
  \mathcal J_2&=\{G,\mathcal H_p(K^o)\},\\
  \mathcal K_2&=\{k^sK^s,(k^s)^2\mathcal H_p(K^o)\}.
 \end{align*}
\end{enumerate}
O verificador confirma as duas provas e compara
\begin{align*}
 k^sZ_?\stackrel{?}{=}(k^s)^2\mathcal H_p(K^o).
\end{align*}
A igualdade indica que a imagem testada gastou a saída; a desigualdade exclui
essa imagem, sob as hipóteses das provas. Para afirmar que a saída continua
não gasta seria necessário repetir o teste para todas as entradas em que
\(K^o\) aparece como membro do anel. Não há comando, formato serializado nem
verificador desta construção no código Monero consultado.

\subsection{Provar a relação entre um endereço e um sub-endereço}
\label{subsec:proofs-address-subaddress-correspond-subaddressproof}

Para um endereço principal com
\(K^v=k^vG\) e um sub-endereço que satisfaz
\begin{align*}
 K^{v,i}=k^vK^{s,i},
\end{align*}
uma prova de igualdade de logaritmos poderia usar
\begin{align*}
 x&=k^v, & J_1&=G, & J_2&=K^{s,i},\\
 P_1&=K^v, & P_2&=K^{v,i}.
\end{align*}
O traslado teria de incluir uma separação de domínio própria, os dois
endereços e a mensagem do verificador. Uma prova válida ligaria o
sub-endereço à chave de vista principal sem divulgar \(k^v\). Esta
\texttt{SubaddressProof} não está implementada no código consultado.

\subsection{Limite de uma auditoria completa}
\label{subsec:audit-framework}
\label{sec:proofs-monero-audit-framework}

Com provas de relação de sub-endereços e provas de não-gasto implementadas,
um provador poderia apresentar uma auditoria mais selectiva:

\begin{enumerate}
 \item declarar os endereços principais e sub-endereços abrangidos e provar
 as relações entre eles;

 \item usar provas de entrada, ou entregar chaves privadas de vista, para
 identificar as saídas recebidas;

 \item testar as imagens de chave de todas as entradas em cujos anéis essas
 saídas aparecem;

 \item fornecer provas de saída para os destinatários cuja divulgação tenha
 sido autorizada.
\end{enumerate}

Este procedimento não está disponível como sistema na carteira. Mesmo que
fosse implementado, não provaria que a lista de contas declarada é completa:
um provador pode controlar outras sementes, endereços ou carteiras. A
completude exige controlos externos ao protocolo. A entrega de uma chave
privada de vista simplifica parte da análise, mas permite observar todas as
recepções passadas e futuras da conta e, sem imagens de chave importadas, não
determina sozinha quais saídas foram gastas.
